\subsection{One comparison rule for a finite set of series}
\label{sec:ordered-eps-method}

The preference shares failed because the options and debt had already reduced
continuing EPS. An implementation needs to reproduce that changing comparison
for any finite number of series. Assume that attribution and instrument-specific
calculations have supplied each series' effects on the same reporting
population, ordinary class, currency scale and retrospective share basis.
Assume also that these effects are fixed and additive: admitting one series
does not change another's conversion terms or earnings consequences. These
assumptions fit the three series above; participating rights or mutually
exclusive conversion choices require further analysis before this procedure.

Write the current continuing ordinary earnings as $N$ and weighted shares
as $D>0$. For one candidate, let $a$ be its signed continuing-earnings
increment and $b>0$ its already time-weighted share increment. Both the
current and enlarged denominators are positive. A strict decrease in EPS is
therefore equivalent to
\begin{align*}
 \frac{N+a}{D+b}<\frac ND
 &\quad\Longleftrightarrow\quad D(N+a)<N(D+b)\\
 &\quad\Longleftrightarrow\quad aD<Nb
 \quad\Longleftrightarrow\quad \frac ab<\frac ND.
\end{align*}
The cancellation uses $DN=ND$ and does not divide by earnings, which may be
zero or negative. It consequently handles a loss without a separate sign
shortcut. For example, if $N=-20$, $D=10$, $a=-15$ and $b=5$, the candidate
changes EPS from $-2$ to $-35/15=-7/3$: it increases the loss per share and
is dilutive. A zero earnings increment with the same positive share increment
would reduce that loss and be excluded. Equality in the comparison is neutral
and is also excluded under the strict dilution rule
\citep[IAS~33, paragraphs 41--44]{eu20231803}.

Sort candidates by their exact ratio $a/b$ in ascending order. Break equal
ratios by a stable series identifier, then test each against the current
$N,D$. On admission, replace $N,D$ with $N+a,D+b$ before proceeding.
Use rational arithmetic for this decision; rounding an incremental ratio to
the displayed precision can change an equality into an apparent decrease.
The debt in the example has the exact ratio $365/184$, and the preference
shares are compared with $7665/1684$, not with the earlier basic ratio.
The share increment already contains its time weight; the selector must not
apply the issue-date fraction a second time.

Why does this order suffice under the additive assumptions? Admission makes
the new ratio a weighted average of the old ratio and $a/b$, strictly between
them. The new ratio remains above the admitted ratio. Every earlier admitted
ratio is therefore below the final ratio, say $q$. At the first exclusion,
the next ratio is at least $q$; every later ratio is at least as large and
is excluded too. If $S$ is the selected set and $N_0,D_0$ the basic amounts,
\[
 N_0-qD_0+\sum_{i\in S}(a_i-qb_i)=0.
\]
Each selected summand is negative; every unselected summand is nonnegative.
Replacing $S$ by any other subset can only increase the left side. Dividing
by that subset's positive denominator shows that its continuing EPS is at
least $q$. This is a local derivation for fixed independent effects, not a
claim that arbitrary instruments satisfy those assumptions. Neutral series
can create other subsets with the same minimum; strict admission excludes them.

Total earnings have their own supplied increment for each series. Carry
those increments through the same selected population, then subtract
continuing earnings to obtain discontinued earnings. Selecting again on
total profit would contradict the continuing-operations control just
illustrated \citep[IAS~33, paragraph A3]{eu20231803}.

\subsection{Keep unresolved instruments in the report}
\label{sec:ordered-eps-report}

The proposed \code{OrderedEPS/1} report accepts attributed basic earnings,
weighted ordinary shares and a separately supported complete instrument
inventory. Each item remains visible as a candidate, a supported no-effect
item, an unresolved item, or an item requiring an unsupported method.
A no-effect item has zero earnings and share increments; it is reported
without dividing by zero. A zero-share item with a nonzero earnings effect
requires another method. An empty inventory also requires its own completeness
finding. These findings are inputs to the calculation: the report cannot
discover omitted instruments from a list supplied as complete.

The request fixes the entity, book and dated head, evidence view, reporting
population and period, ordinary class, share basis, currency and scale, and
historical rule version. Basic and instrument records repeat a hash of that
context. Every premise also points to exact source bytes and a source anchor.
Matching hashes establish which evidence was used; interpretation, completeness
and operative authority still require their own findings. Potential-share
windows must lie within the period and must not overlap. A converted series
also identifies its subsequent actual-share window and confirms that those
actual shares already enter basic EPS, as in the October conversion above.
The date checks prevent overlap; instrument methods must still establish the
amount and weighting of each increment.

A completed conditional report returns basic and diluted continuing, total
and discontinued EPS, their numerators and common denominators, and the
ordered comparison for every series. Each comparison retains its increments,
the amounts before the test, the amounts if admitted, and the admission
decision. Missing attribution can prevent both basic and diluted EPS. Missing
instrument consequences or an incomplete inventory preserve an independently
ready basic result but leave diluted EPS unresolved. Known dependent effects
or participation terms produce an unsupported result; both kinds of obstacle
are reported together when present. Invalid bindings, malformed fractions,
duplicate or omitted inventoried series, and overlapping windows reject the
request. None of these outcomes returns a partial selection as complete.

The exact fields and result states are specified in
Section~\ref{sec:ordered-eps-records}. This is a read-only proposed report:
it creates no journal or share and changes no book. Its finite reference
checks cover the supplied arithmetic, ordering and refusal behavior.
Applying it to a company still requires instrument-specific preprocessing,
complete facts and applicable authority; joining it to a company's dated
earnings, tax and share records is a further implementation obligation.
